\documentclass[10pt,a4paper]{amsart}
\usepackage{amsmath,amssymb,amsthm,graphicx}
\usepackage{verbatim}
\newenvironment{reference}{\comment}{\endcomment}
%\newenvironment{reference}{}{}
\newenvironment{slogan}{\comment}{\endcomment}
\newenvironment{history}{\comment}{\endcomment}
% For commutative diagrams we use Xy-pic
\usepackage[all]{xy}

% We use 2cell for 2-commutative diagrams.
\xyoption{2cell}
% Theorem environments.
%
\theoremstyle{plain}
\newtheorem{theorem}[subsection]{Theorem}
\newtheorem{proposition}[subsection]{Proposition}
\newtheorem{lemma}[subsection]{Lemma}

\theoremstyle{definition}
\newtheorem{definition}[subsection]{Definition}
\newtheorem{example}[subsection]{Example}
\newtheorem{exercise}[subsection]{Exercise}
\newtheorem{situation}[subsection]{Situation}

\theoremstyle{remark}
\newtheorem{remark}[subsection]{Remark}
\newtheorem{remarks}[subsection]{Remarks}

\numberwithin{equation}{subsection}

% Macros
%
\def\lim{\mathop{\mathrm{lim}}\nolimits}
\def\colim{\mathop{\mathrm{colim}}\nolimits}
\def\Spec{\mathop{\mathrm{Spec}}}
\def\Hom{\mathop{\mathrm{Hom}}\nolimits}
\def\Ext{\mathop{\mathrm{Ext}}\nolimits}
\def\SheafHom{\mathop{\mathcal{H}\!\mathit{om}}\nolimits}
\def\SheafExt{\mathop{\mathcal{E}\!\mathit{xt}}\nolimits}
\def\Sch{\mathit{Sch}}
\def\Mor{\mathop{\mathrm{Mor}}\nolimits}
\def\Ob{\mathop{\mathrm{Ob}}\nolimits}
\def\Sh{\mathop{\mathit{Sh}}\nolimits}
\def\NL{\mathop{N\!L}\nolimits}
\def\CH{\mathop{\mathrm{CH}}\nolimits}
\def\proetale{{pro\text{-}\acute{e}tale}}
\def\etale{{\acute{e}tale}}
\def\QCoh{\mathit{QCoh}}
\def\Ker{\mathop{\mathrm{Ker}}}
\def\Im{\mathop{\mathrm{Im}}}
\def\Coker{\mathop{\mathrm{Coker}}}
\def\Coim{\mathop{\mathrm{Coim}}}

% Boxtimes
%
\DeclareMathSymbol{\boxtimes}{\mathbin}{AMSa}{"02}

%
% Macros for moduli stacks/spaces
%
\def\QCohstack{\mathcal{QC}\!\mathit{oh}}
\def\Cohstack{\mathcal{C}\!\mathit{oh}}
\def\Spacesstack{\mathcal{S}\!\mathit{paces}}
\def\Quotfunctor{\mathrm{Quot}}
\def\Hilbfunctor{\mathrm{Hilb}}
\def\Curvesstack{\mathcal{C}\!\mathit{urves}}
\def\Polarizedstack{\mathcal{P}\!\mathit{olarized}}
\def\Complexesstack{\mathcal{C}\!\mathit{omplexes}}
% \Pic is the operator that assigns to X its picard group, usage \Pic(X)
% \Picardstack_{X/B} denotes the Picard stack of X over B
% \Picardfunctor_{X/B} denotes the Picard functor of X over B
\def\Pic{\mathop{\mathrm{Pic}}\nolimits}
\def\Picardstack{\mathcal{P}\!\mathit{ic}}
\def\Picardfunctor{\mathrm{Pic}}
\def\Deformationcategory{\mathcal{D}\!\mathit{ef}}
\usepackage{cleveref}
\crefname{originallemma}{original lemma}{original lemmas}
\Crefname{originallemma}{Original lemma}{Original lemmas}
\makeatletter
\newlabel{lemma-isom-locally-finite-type}{{lemma-isom-locally-finite-type}{source}{}{sourcepointer-1}{}}
\newlabel{lemma-isom-locally-finite-type@cref}{{[originallemma][0][]lemma-isom-locally-finite-type}{[1][0][]source}{}{}{}}
\makeatother
\crefname{originalreference}{original reference}{original references}
\Crefname{originalreference}{Original reference}{Original references}
\makeatletter
\newlabel{spaces-morphisms-lemma-base-change-finite-type}{{spaces-morphisms-lemma-base-change-finite-type}{source}{}{sourcepointer-2}{}}
\newlabel{spaces-morphisms-lemma-base-change-finite-type@cref}{{[originalreference][0][]spaces-morphisms-lemma-base-change-finite-type}{[1][0][]source}{}{}{}}
\makeatother
\crefname{originalreference}{original reference}{original references}
\Crefname{originalreference}{Original reference}{Original references}
\makeatletter
\newlabel{spaces-morphisms-lemma-permanence-finite-type}{{spaces-morphisms-lemma-permanence-finite-type}{source}{}{sourcepointer-3}{}}
\newlabel{spaces-morphisms-lemma-permanence-finite-type@cref}{{[originalreference][0][]spaces-morphisms-lemma-permanence-finite-type}{[1][0][]source}{}{}{}}
\makeatother
\title{The diagonal of a morphism of algebraic stacks is representable by algebraic spaces and locally of finite type}
\author{The Stacks Project Authors}
\date{stacks-project:stacks-morphisms.tex:local-snapshot-20261012; GNU FDL 1.2 or later}
\begin{document}
\maketitle
\noindent\textit{Editorial note.} All delivered mathematical text is the author's own, copied from the pinned Stacks Project chapter source sources/stacks-algebraic-stacks-morphisms/pinned-original-source.tex; no statement or proof step was rewritten, reordered or removed.\par\medskip
\noindent\textit{Editorial note.} Preamble provenance: the Stacks Project keeps one single shared preamble for every chapter; this card takes it directly from the staged sources/stacks-algebraic-stacks-morphisms/preamble.tex as three verbatim spans (lines 8-12, 14-18 and 154-217), so no derived file was written and nothing was added or redefined. The wrapper class and the base packages are supplied by the assembler.\par\medskip
\noindent\textit{Editorial note.} Bibliography provenance: neither the counted statement nor its proof contains a citation (this lane rescans the whole span before writing the card), so the card carries no bibliography and no entry was reconstructed.\par\medskip
\noindent\textit{Editorial note.} Selection provenance: this card was opened under the file-based control contract (CONTROL.md), with the mechanical pre-card gate run on its statement and proof interval before the card was written.\par\medskip
\noindent\textit{Allowed standard prerequisites (editorial).}\begin{itemize}
\item The statement that the isomorphism sheaf of two objects of an algebraic stack is an algebraic space. Used inside the counted proof (source lines 365-369). Makes the two isomorphism sheaves appearing in the fibre square algebraic spaces.
\item The statement that the isomorphism space of two objects is locally of finite type. Invoked inside the counted proof (source lines 391-393). Gives local finite type of the right vertical arrow of the displayed square.
\item The base change and composition permanence statements for locally of finite type morphisms of algebraic spaces. Invoked at the end of the counted proof (source lines 394-398) as cross-chapter pointers. Transports local finite type from the right vertical arrow to the left one.
\end{itemize}\medskip
\ifdefined\hypertarget\hypertarget{sourcepointer-1}{}\fi
\noindent\textit{Source reference.} Original lemma lemma-isom-locally-finite-type: sources/stacks-algebraic-stacks-morphisms/pinned-original-source.tex, source lines 278--278. Navigation pointer only: the counted text refers to the Stacks item lemma-isom-locally-finite-type, which is an existing labeled result of this same chapter and is not part of this card. The pointer carries the target's own label as its number and its definition line as its span; its content is not reproduced here.\par\medskip
\ifdefined\hypertarget\hypertarget{sourcepointer-2}{}\fi
\noindent\textit{Source reference.} Original reference spaces-morphisms-lemma-base-change-finite-type: sources/stacks-algebraic-stacks-morphisms/pinned-original-source.tex, source lines 401--401. Navigation pointer only: the counted text refers to the Stacks item spaces-morphisms-lemma-base-change-finite-type, which belongs to another chapter of the project and is not part of this card. The pointer carries the target's own label as its number because this chapter's staged snapshot does not contain the other chapter's tag table; its content is not reproduced here.\par\medskip
\ifdefined\hypertarget\hypertarget{sourcepointer-3}{}\fi
\noindent\textit{Source reference.} Original reference spaces-morphisms-lemma-permanence-finite-type: sources/stacks-algebraic-stacks-morphisms/pinned-original-source.tex, source lines 307--307. Navigation pointer only: the counted text refers to the Stacks item spaces-morphisms-lemma-permanence-finite-type, which belongs to another chapter of the project and is not part of this card. The pointer carries the target's own label as its number because this chapter's staged snapshot does not contain the other chapter's tag table; its content is not reproduced here.\par\medskip
\begin{lemma}
\label{lemma-properties-diagonal}
\begin{slogan}
Diagonals of morphisms of algebraic stacks are representable by
algebraic spaces and locally of finite type.
\end{slogan}
Let $f : \mathcal{X} \to \mathcal{Y}$ be a morphism of algebraic stacks.
Then
\begin{enumerate}
\item $\Delta_f$ is representable by algebraic spaces,
and
\item $\Delta_f$ is locally of finite type.
\end{enumerate}
\end{lemma}

\begin{proof}
Let $T$ be a scheme and let
$a : T \to \mathcal{X} \times_\mathcal{Y} \mathcal{X}$
be a morphism. By definition of the fibre product and the
$2$-Yoneda lemma the morphism $a$ is given by a triple
$a = (x, x', \alpha)$ where $x, x'$ are objects of $\mathcal{X}$
over $T$, and $\alpha : f(x) \to f(x')$ is a morphism in the fibre
category of $\mathcal{Y}$ over $T$. By definition of an algebraic
stack the sheaves $\mathit{Isom}_\mathcal{X}(x, x')$ and
$\mathit{Isom}_\mathcal{Y}(f(x), f(x'))$ are algebraic spaces
over $T$. In this language $\alpha$ defines a section of the morphism
$\mathit{Isom}_\mathcal{Y}(f(x), f(x')) \to T$. A $T'$-valued point of
$\mathcal{X} \times_{\mathcal{X} \times_\mathcal{Y} \mathcal{X}, a} T$
for $T' \to T$ a scheme over $T$ is the same thing as an isomorphism
$x|_{T'} \to x'|_{T'}$ whose image under $f$ is $\alpha|_{T'}$.
Thus we see that
\begin{equation}
\label{equation-diagonal}
\vcenter{
\xymatrix{
\mathcal{X} \times_{\mathcal{X} \times_\mathcal{Y} \mathcal{X}, a} T
\ar[d] \ar[r] &
\mathit{Isom}_\mathcal{X}(x, x') \ar[d] \\
T\ar[r]^-\alpha &
\mathit{Isom}_\mathcal{Y}(f(x), f(x'))
}
}
\end{equation}
is a fibre square of sheaves over $T$. In particular we see that
$\mathcal{X} \times_{\mathcal{X} \times_\mathcal{Y} \mathcal{X}, a} T$
is an algebraic space which proves part (1) of the lemma.

\medskip\noindent
To prove the second statement we have to show that the left
vertical arrow of Diagram (\ref{equation-diagonal}) is locally
of finite type. By
Lemma \ref{lemma-isom-locally-finite-type}
the algebraic space $\mathit{Isom}_\mathcal{X}(x, x')$
is locally of finite type over $T$. Hence the right vertical arrow of
Diagram (\ref{equation-diagonal}) is locally of finite type, see
Morphisms of Spaces, Lemma \ref{spaces-morphisms-lemma-permanence-finite-type}.
We conclude by
Morphisms of Spaces,
Lemma \ref{spaces-morphisms-lemma-base-change-finite-type}.
\end{proof}

\end{document}
